% lecture_7.tex

\section*{Cursul VII}

\begin{corolar}{Lema Banach}
    $$ \forall F \in \M_n(\R) : \|F\| < 1, $$
    $$ \|(I_n - F)^{-1} - I_n\| \leq \frac{\|F\|}{1 - \|F\|} $$
\end{corolar}

\begin{demonstratie}
    \begin{align*}
        (I_n - F)^{-1} &= \sum_{k=0}^{\infty} F^k  \\
        (I_n - F)^{-1} - I_n &= \sum_{k=1}^{\infty} F^k\\
        &= F \sum_{k=0}^{\infty} F^k \\
        &= F (I_n - F)^{-1}.
    \end{align*}

    Atunci:
    \begin{align*}
        \|(I_n - F)^{-1} - I_n\| &= \| F (I_n - F)^{-1} \| \\
        &\leq \| F \| \| (I_n - F)^{-1} \| \\
        &\leq \| F \| \frac{1}{1 - \|F\|} \\
        &= \frac{\|F\|}{1 - \|F\|}.
    \end{align*}
\end{demonstratie}

\begin{teorema}{1}

Daca $A \in \M_n(\R)$ este inversabila si $E \in \M_n(\R)$ a.i. $r := \|A^{-1}E\| < 1$, atunci
\begin{enumerate}
    \item[(i)] $A+E \in \M_n(\R)$ este inversabila
    \item[(ii)] $ \| (A+E)^{-1} - A^{-1} \| \leq \frac{\|E\| \cdot \|A^{-1}\|^2}{1-r} $
\end{enumerate}
\end{teorema}

\begin{demonstratie}

    \textbf{i)} P.p. ca $ A + E \rightarrow $ singulara. Deci $ (\exists) x \in \R^{n} \neq 0 \quad a.i. \quad (A+E)x = 0 $.
    \begin{align*}
        Ax + Ex &= 0 \\
        Ax &= -Ex \\
        x &= -A^{-1}Ex \\
        \|x\| &= \|-A^{-1}Ex\| \\
        \|x\| &\leq \|A^{-1}E\| \|x\| \\ 
        1 &\leq \|A^{-1}E\|, 
    \end{align*}
    ceea ce e o contradictie $ \implies A $ este nesingulara si inversabila.

    \bigskip

    \textbf{ii)}
    $$ A + E = A(I_n + A^{-1}E) = A[I_n - (-A^{-1}E)] $$.

    Considerand $F \in \M_n(\R)$, $F = -A^{-1}E$, rezulta ca:
    \begin{align*}
        \|(A+E)^{-1} - A^{-1}\| &= \| [A(I_n - F)]^{-1} - A^{-1} \| \\
        &= \| (I_n - F)^{-1} \cdot A^{-1} - A^{-1} \| \\
        &= \| A^{-1} [(I_n - F)^{-1} - I_n] \| \\
        &\leq \| A^{-1} \| \| (I_n - F)^{-1} - I_n \|,
    \end{align*}
    care, conform corolarului Lemei Banach,
    $$ \leq \|A^{-1}\| \cdot \frac{\|F\|}{1 - \|F\|}. $$

    Inlocuind $F$ cu $-A^{-1}E$ si tinand cont de invarianta normei matriceale la schimbarea de semn,

    \begin{align*}
        \|(A+E)^{-1} - A^{-1}\| &\leq \|A^{-1}\| \cdot \frac{\|A^{-1}E\|}{1 - \|A^{-1}E\|} \\
        &\leq \frac{\|E\| \|A^{-1}\|^2}{1 - r}.
    \end{align*}

\end{demonstratie}

\begin{exercitiu}{I} 
    Sa se determine $h\p{k}$ dat de metoda Newton-Jacobi cu $ m \in \N $ pasi, i.e., Jacobi($m$).
\end{exercitiu}

\begin{demonstratie}

    Conform cursului, pentru metoda Newton-Richardson Stationara:
    $$ h\p{k} := - \sum_{j=0}^{m-1} (M\p{k})^j \cdot (P\p{k})^{-1} \cdot F(x\p{k}) $$

    In particular, pentru metoda Newton-Jacobi:
    \begin{align*}
        h\p{k} &:= - \sum_{j=0}^{m-1} (M_J\p{k})^j \cdot (P_J\p{k})^{-1} \cdot F(x\p{k}) \\
        &= \sum_{j=0}^{m-1} \left[ (D\p{k})^{-1} (L\p{k} + U\p{k}) \right]^j \cdot (D\p{k})^{-1} \cdot F(x\p{k}),
    \end{align*}
    unde $ D $, $ U $ si $ L $ sunt partile aditive ale matricii $ J_F (x\p{k})$.

\end{demonstratie}